
\documentclass{article}
\usepackage{colortbl}
\usepackage{makecell}
\usepackage{multirow}
\usepackage{supertabular}

\begin{document}

\newcounter{utterance}

\centering \large Interaction Transcript for game `ifeval', experiment `ifeval', episode 37 with selfrefine{-}llama{-}3.1{-}8b{-}instruct.
\vspace{24pt}

{ \footnotesize  \setcounter{utterance}{1}
\setlength{\tabcolsep}{0pt}
\begin{supertabular}{c@{$\;$}|p{.15\linewidth}@{}p{.15\linewidth}p{.15\linewidth}p{.15\linewidth}p{.15\linewidth}p{.15\linewidth}}
    \# & \multicolumn{2}{c}{Player} && \multicolumn{2}{c}{Game Master} \\
    \hline

    \theutterance \stepcounter{utterance}  
    & & & \multicolumn{4}{p{0.6\linewidth}}{
        \cellcolor[rgb]{0.9,0.9,0.9}{
            \makecell[{{p{\linewidth}}}]{
                \texttt{\tiny{[P1$\langle$GM]}}
                \texttt{What sentiments existed in the Croatian and Serbian nations towards the Austrian Empire at the end of the 19th century? Make sure to wrap your entire response with double quotes and use at least 800 words.} \\
            }
        }
    }
     \\ \\

    \theutterance \stepcounter{utterance}  
    & \multicolumn{4}{p{0.6\linewidth}}{
        \cellcolor[rgb]{0.561,0.737,0.561}{
            \makecell[{{p{\linewidth}}}]{
                \texttt{\tiny{[P1$\rangle$GM]}}
                \texttt{"Throughout the 19th century, the Croatian and Serbian nations experienced a complex and often tumultuous relationship with the Austrian Empire. The sentiments towards the empire varied greatly among the two nations, reflecting their distinct historical, cultural, and political contexts.} \\
\\ 
\texttt{In Croatia, the sentiment towards the Austrian Empire was largely one of frustration and disillusionment. The Croatian people had long been seeking greater autonomy and self{-}governance within the empire, but their efforts were consistently met with resistance and repression from the Austrian authorities. The Croatian National Revival, which emerged in the early 19th century, sought to promote Croatian language, culture, and identity, but it was often seen as a threat to the empire's unity and authority.} \\
\\ 
\texttt{In contrast, the Serbian people had a more complex and nuanced relationship with the Austrian Empire. On the one hand, the Serbian people had long been seeking greater autonomy and self{-}governance within the empire, but they were also deeply concerned about the empire's policies towards the Balkans. The Serbian people saw the Austrian Empire as a potential protector against the Ottoman Empire, which had long been a threat to their national security and identity.} \\
\\ 
\texttt{However, the Serbian people were also deeply concerned about the empire's treatment of the Croatian people, whom they saw as a rival national group. The Serbian people felt that the Croatian people were being favored by the empire, and that their own national interests were being ignored. This led to a deep{-}seated mistrust and hostility between the two nations, which would ultimately contribute to the outbreak of the Balkan Wars and the dissolution of the empire.} \\
\\ 
\texttt{The Austrian Empire's policies towards the Balkans were a major source of tension between the Croatian and Serbian people. The empire's policy of "divide and rule" sought to maintain control over the region by playing off the different national groups against each other. This led to a complex web of alliances and rivalries between the different national groups, which ultimately contributed to the outbreak of conflict.} \\
\\ 
\texttt{The Croatian people were also concerned about the empire's treatment of the Hungarian people, whom they saw as a rival national group. However, the Serbian people were more focused on their national identity and interests, and they were less concerned about the empire's policies towards the Hungarian people.} \\
\\ 
\texttt{The empire's policies towards the Balkans ultimately contributed to the outbreak of the Balkan Wars and the dissolution of the empire. The complex web of alliances and rivalries between the different national groups, which was fueled by the empire's policy of "divide and rule," ultimately led to the outbreak of conflict.} \\
\\ 
\texttt{In the late 19th century, the Croatian and Serbian people were both seeking greater autonomy and self{-}governance within the empire. However, the empire's policies towards the Balkans were ultimately designed to maintain control over the region, rather than to promote the interests of the different national groups.} \\
\\ 
\texttt{The Croatian people's sentiment towards the Austrian Empire was characterized by frustration and disillusionment, while the Serbian people's sentiment was more complex and nuanced. The empire's policies towards the Balkans were a major source of tension between the two nations, and ultimately contributed to the outbreak of the Balkan Wars and the dissolution of the empire."} \\
            }
        }
    }
    & & \\ \\

    \theutterance \stepcounter{utterance}  
    & & & \multicolumn{2}{p{0.3\linewidth}}{
        \cellcolor[rgb]{0.9,0.9,0.9}{
            \makecell[{{p{\linewidth}}}]{
                \texttt{\tiny{[GM$|$GM]}}
                {'startend:quotation': {}, 'length_constraints:number_words': {'relation': 'at least', 'num_words': 800}}
            }
        }
    }
    & & \\ \\

    \theutterance \stepcounter{utterance}  
    & & & \multicolumn{2}{p{0.3\linewidth}}{
        \cellcolor[rgb]{0.9,0.9,0.9}{
            \makecell[{{p{\linewidth}}}]{
                \texttt{\tiny{[GM$|$GM]}}
                \texttt{game\_result = LOSE} \\
            }
        }
    }
    & & \\ \\

\end{supertabular}
}

\end{document}
